\documentclass[11pt]{article}
\usepackage[utf8]{inputenc}
\usepackage[spanish]{babel}
\usepackage{amsmath,amsthm,amsfonts,amssymb,amscd}
\usepackage{multirow,booktabs}
\usepackage[table]{xcolor}
\usepackage{fullpage}
\usepackage{lastpage}
\usepackage{enumitem}
\usepackage{fancyhdr}
\usepackage{mathrsfs}
\usepackage{wrapfig}
\usepackage{setspace}
\usepackage{calc}
\usepackage{multicol}
\usepackage[retainorgcmds]{IEEEtrantools}
\usepackage[margin=3cm]{geometry}
\usepackage{empheq}
\usepackage{framed}
\usepackage[most]{tcolorbox}
\usepackage{soul}
\usepackage{longtable}
\usepackage{array}
\usepackage{ragged2e}
\newlength{\tabcont}
\setlength{\parindent}{0.0in}
\setlength{\parskip}{0.05in}
\colorlet{shadecolor}{orange!15}
\parindent 0in
\parskip 12pt
\geometry{margin=1in, headsep=0.25in}

\newcolumntype{L}[1]{>{\RaggedRight\arraybackslash}p{#1}}

\begin{document}
\setcounter{section}{0}

\thispagestyle{empty}

\begin{center}
{\LARGE \bf Ruta Crítica de Ejecución}\\[4pt]
{\large Tesis I}\\[2pt]
Roberto Treviño Cervantes
\end{center}

\section{Contexto}

La presente ruta crítica organiza las actividades de tesis para el periodo comprendido entre mayo y septiembre de 2026, cubriendo los cinco meses restantes antes de la entrega del manuscrito. Cada hito se asocia a un mes específico y representa un entregable verificable que alimenta al siguiente. El alcance es exclusivamente computacional: refinamiento de modelo, simulación, validación, análisis y escritura, todos sobre el gemelo digital del escáner de fotolitografía.

\section{Hitos Mensuales}

\begingroup
\small
\setlength{\tabcolsep}{6pt}
\renewcommand{\arraystretch}{1.5}
\begin{longtable}{L{2.5cm} L{3.5cm} L{6.5cm}}
\caption{Ruta crítica de ejecución mayo--septiembre 2026.}
\label{tab:ruta} \\
\toprule
\textbf{Mes} & \textbf{Hito} & \textbf{Descripción del entregable} \\
\midrule
\endfirsthead
\multicolumn{3}{c}{\small\textit{(continuación de la Tabla~\ref{tab:ruta})}} \\[4pt]
\toprule
\textbf{Mes} & \textbf{Hito} & \textbf{Descripción del entregable} \\
\midrule
\endhead
\bottomrule
\endlastfoot

Mayo 2026 &
Refinamiento del modelo &
Ajuste del modelo multi-cuerpo y del generador de perfiles de cuarto orden con 15 segmentos. Verificación de parámetros inerciales, rigideces y amortiguamientos. \\[4pt]

Junio 2026 &
Simulaciones iniciales &
Ejecución de la simulación con barridos sistemáticos de los parámetros cinemáticos $(v_{max},\, a_{max},\, j_{max},\, s_{max})$. Generación del primer conjunto de mapas de error de seguimiento. \\[4pt]

Julio 2026 &
Validación &
Entrenamiento y validación cruzada de la CNN sobre WM-811K como función de costo sustituta. Validación de la correspondencia entre patrones de error de trayectoria y morfologías de defecto. \\[4pt]

Agosto 2026 &
Análisis &
Ejecución del PSO acoplado a la CNN como función de costo. Análisis comparativo de las trayectorias óptimas en términos de \textit{yield} estimado frente a \textit{throughput}. \\[4pt]

Septiembre 2026 &
Escritura &
Redacción del manuscrito final de tesis, integración de figuras, tablas y referencias. \\[4pt]

\end{longtable}
\endgroup

\end{document}
